\documentclass[11pt,letterpaper]{article}
\usepackage[T1]{fontenc}
\usepackage{lmodern}
\usepackage[margin=1in,headheight=14pt]{geometry}
\usepackage{amsmath,amssymb,amsthm,mathtools}
\usepackage{microtype,booktabs,array,aliascnt}
\usepackage{xcolor}
\usepackage{fancyhdr}
\usepackage{enumitem}
\usepackage[colorlinks=true,linkcolor=blue!45!black,citecolor=blue!45!black,urlcolor=blue!45!black]{hyperref}
\usepackage[nameinlink,noabbrev]{cleveref}
\hypersetup{pdftitle={Apery Numbers under Exponentiation and Series Reversion},pdfauthor={ChatGPT},pdfsubject={Proofs of two OEIS A267220 supercongruence conjectures}}
\pagestyle{fancy}
\fancyhf{}
\fancyhead[L]{\small Ap\'ery numbers and two OEIS conjectures}
\fancyhead[R]{\small\thepage}
\renewcommand{\headrulewidth}{0.3pt}
\setlength{\parskip}{3pt}
\setlength{\emergencystretch}{2em}
\numberwithin{equation}{section}
\newtheorem{theorem}{Theorem}[section]
\newaliascnt{proposition}{theorem}
\newtheorem{proposition}[proposition]{Proposition}
\aliascntresetthe{proposition}
\crefname{proposition}{proposition}{propositions}
\Crefname{proposition}{Proposition}{Propositions}
\newaliascnt{lemma}{theorem}
\newtheorem{lemma}[lemma]{Lemma}
\aliascntresetthe{lemma}
\crefname{lemma}{lemma}{lemmas}
\Crefname{lemma}{Lemma}{Lemmas}
\newaliascnt{corollary}{theorem}
\newtheorem{corollary}[corollary]{Corollary}
\aliascntresetthe{corollary}
\crefname{corollary}{corollary}{corollaries}
\Crefname{corollary}{Corollary}{Corollaries}
\theoremstyle{definition}
\newaliascnt{definition}{theorem}
\newtheorem{definition}[definition]{Definition}
\aliascntresetthe{definition}
\crefname{definition}{definition}{definitions}
\Crefname{definition}{Definition}{Definitions}
\newaliascnt{example}{theorem}
\newtheorem{example}[example]{Example}
\aliascntresetthe{example}
\crefname{example}{example}{examples}
\Crefname{example}{Example}{Examples}
\theoremstyle{remark}
\newaliascnt{remark}{theorem}
\newtheorem{remark}[remark]{Remark}
\aliascntresetthe{remark}
\crefname{remark}{remark}{remarks}
\Crefname{remark}{Remark}{Remarks}
\newcommand{\Z}{\mathbb Z}
\newcommand{\Q}{\mathbb Q}
\newcommand{\Zp}{\mathbb Z_p}
\newcommand{\Qp}{\mathbb Q_p}
\newcommand{\vp}{\operatorname{ord}_p}
\newcommand{\vtwo}{\operatorname{ord}_2}
\newcommand{\Rev}{\operatorname{Rev}}
\newcommand{\Li}{\operatorname{Li}}
\newcommand{\coeff}[2]{[#1^{#2}]}
\newcommand{\dd}{\mathop{}\!d}
\title{\textbf{Ap\'ery Numbers under Exponentiation\texorpdfstring{\\}{ }and Series Reversion}\\[8pt]
\large Proofs of two OEIS supercongruence conjectures\\
\normalsize A self-contained framing argument for A267220}
\author{ChatGPT}
\date{20 September 2026}
\begin{document}
\maketitle
\begin{abstract}
Let $a_n=\sum_{k=0}^n\binom nk^2\binom{n+k}k^2$ be the classical Ap\'ery numbers,
let $G(z)=\exp(\sum_{n\ge1}a_nz^n/n)$, and put
$F(w)=w/\Rev(zG(z))(w)$. We prove both conjectures recorded by Peter Bala
in OEIS A267220 on 17 October 2024: for every integer $m$, the sequences
$u_m(n)=[z^n]G(z)^{mn}$ and $v_m(n)=[w^n]F(w)^{mn}$ satisfy
order-two prime-power supercongruences. The proof extends both conjectures
to $p=3$, and extends the second to $p=2$ as well. It also gives
parameter-dependent divisibility and proves integrality of the normalized
family $B_t(n)=t^{-1}[z^n]G(z)^{tn}$, with $B_0(n)=a_n$.
An elementary binomial-product argument supplies the initial congruences.
An explicit Frobenius-defect identity then establishes the required
change-of-variable theorem, a rational-coefficient instance of the known
framing principle for $2$-functions. The special behavior at $2$ is
tracked without suppressing a denominator. Small exact examples show
that a uniform cubic modulus is false, and that the first family cannot
be extended unchanged to $2$. A standard-library Python program provides
independent finite checks using integer and rational arithmetic.
\end{abstract}
\vspace{4pt}
\noindent\textbf{Scope and status.} This article supplies mathematical proofs,
not an inference from numerical experiments. The general framing principle
is prior work, not a novelty claim. The specific OEIS applications and
refinements are proved explicitly here. This is an unrefereed research
exposition; no proof-assistant verification or OEIS submission is claimed.
\clearpage
\tableofcontents
\clearpage

\section{The problem and its mathematical setting}\label{sec:problem}
The Ap\'ery numbers associated with $\zeta(3)$ are
\begin{equation}\label{eq:apery}
 a_n=\sum_{k=0}^n\binom nk^2\binom{n+k}k^2
 \qquad(n\ge0).
\end{equation}
They begin $1,5,73,1445,33001,819005,21460825,\ldots$ and form
OEIS A005259 \cite{oeis-apery}. Their role in Ap\'ery's irrationality proof
and their classical supercongruences make them a natural, substantial
example rather than an artificially chosen sequence; see \cite{liu}
for the arithmetic background.

Define the formal power series
\begin{equation}\label{eq:basic-series}
 L(z)=\sum_{n\ge1}\frac{a_n}{n}z^n,
 \qquad G(z)=\exp L(z),
 \qquad W(z)=\sum_{n\ge1}\frac{a_n}{n^2}z^n.
\end{equation}
We use $G$ for the generating function called $A$ in OEIS A267220, to
avoid confusing it with the Ap\'ery numbers themselves. If $h(w)$ is the
unique compositional inverse of $zG(z)$, set
\begin{equation}\label{eq:F}
 h(w)=\Rev(zG(z))(w),\qquad F(w)=\frac{w}{h(w)}.
\end{equation}
Here reversion means composition, not the multiplicative reciprocal.
The first coefficients are
\begin{align*}
 G(z)&=1+5z+49z^2+685z^3+11807z^4+232771z^5+\cdots,\\
 F(w)&=1+5w+24w^2+200w^3+2430w^4+36096w^5+\cdots.
\end{align*}
In the entry's formula section, Peter Bala's contribution dated
17 October 2024 states the following two conjectures \cite{oeis-target}.
For each integer $m$, let
\begin{equation}\label{eq:uv}
 u_m(n)=[z^n]G(z)^{mn},\qquad
 v_m(n)=[w^n]F(w)^{mn}.
\end{equation}
For every prime $p\ge5$ and positive integers $n,r$, the proposed formulas are
\begin{align}
 u_m(np^r)&\equiv u_m(np^{r-1})\pmod{p^{2r}},\label{eq:conj-u}\\
 v_m(np^r)&\equiv v_m(np^{r-1})\pmod{p^{2r}}.\label{eq:conj-v}
\end{align}
Both remain labeled conjectures in the entry accessed on
20 September 2026. These are the targets of the present article. An older,
different integrality comment in that entry already cites a proof; it is
not being presented here as an open problem.

The exponent in \eqref{eq:uv} grows with the index, so these are not merely
fixed powers of one integral series. Reversion adds a nonlinear operation.
Furthermore, the original parameter range includes negative integers,
$m=0$, and parameters divisible by the modulus prime. A proof by division
by $m$ or $m-1$ in a finite field would miss some of those cases.

The appropriate organizing object is a \emph{$2$-function}, a formal
series whose coefficients obey a prime-by-prime divisibility condition.
Schwarz, Vologodsky, and Walcher established preservation of this
structure under framing, a change of variable involving compositional
inversion \cite{svw2013,svw2017}. Our proof develops exactly the local
rational-coefficient calculation needed here. In particular, it keeps the
unsigned convention dictated by \eqref{eq:uv} and treats its behavior at
$2$ directly.

\section{Main results and conventions}\label{sec:main}
All series manipulations take place formally, initially over $\Q$.
No analytic convergence is assumed. We write $\vp(q)$ for the exponent
of a prime $p$ in a nonzero rational number $q$ and use
$\vp(0)=+\infty$. Thus $q\in\Zp$ means $\vp(q)\ge0$; this is simply a
convenient way of saying that the reduced denominator of $q$ is not
multiple of $p$. A congruence between integers has its usual meaning.
For a variable $z$ put $\delta_z=z\dd/\dd z$.

For an integer $t$ and $n\ge1$, define
\begin{equation}\label{eq:Bt}
 B_t(n)=
 \begin{cases}
 \displaystyle\frac1t[z^n]G(z)^{tn},&t\ne0,\\[5pt]
 a_n,&t=0.
 \end{cases}
\end{equation}
The value at zero is the removable value: the coefficient of $t$ in
$[z^n]\exp(tnL(z))$ is $a_n$. More generally $B_t(n)$ is a rational
polynomial in $t$ of degree at most $n-1$.

\begin{theorem}[Normalized supercongruences]\label{thm:normalized}
For every integer $t$ and every $n\ge1$, $B_t(n)$ is an integer.
For an odd prime $p$ and an integer $N\ge1$ divisible by $p$,
\begin{equation}\label{eq:B-odd}
 B_t(N)\equiv B_t(N/p)\pmod{p^{2\vp(N)}}.
\end{equation}
If $N$ is even, then
\begin{equation}\label{eq:B-two}
 \vtwo\bigl(B_t(N)-B_t(N/2)\bigr)\ge
 \begin{cases}
 2\vtwo(N),&t\text{ even},\\
 2\vtwo(N)-1,&t\text{ odd}.
 \end{cases}
\end{equation}
\end{theorem}

\begin{theorem}[Both OEIS conjectures, with refinements]\label{thm:OEIS}
For all integers $m$ and positive integers $n$,
\begin{equation}\label{eq:identity-uv}
 u_m(n)=mB_m(n),\qquad v_m(n)=mB_{m-1}(n).
\end{equation}
Consequently, \eqref{eq:conj-u} and \eqref{eq:conj-v} hold for every odd
prime, including $3$. In fact, if $m\ne0$, $p$ is odd, and $p\mid N$, then
\begin{align}
 \vp\bigl(u_m(N)-u_m(N/p)\bigr)&\ge2\vp(N)+\vp(m),\label{eq:strong-u}\\
 \vp\bigl(v_m(N)-v_m(N/p)\bigr)&\ge2\vp(N)+\vp(m).\label{eq:strong-v}
\end{align}
The second OEIS congruence also holds for $p=2$, for every integer $m$.
More precisely, for nonzero $m$ and even $N$,
\begin{align}
 \vtwo\bigl(u_m(N)-u_m(N/2)\bigr)&\ge
 \begin{cases}
 2\vtwo(N)+\vtwo(m),&m\text{ even},\\
 2\vtwo(N)-1,&m\text{ odd},
 \end{cases}\label{eq:u-two}\\
 \vtwo\bigl(v_m(N)-v_m(N/2)\bigr)&\ge
 \begin{cases}
 2\vtwo(N)+\vtwo(m)-1,&m\text{ even},\\
 2\vtwo(N),&m\text{ odd}.
 \end{cases}\label{eq:v-two}
\end{align}
For $m=0$, both sequences are zero at positive indices, so all these
congruences are automatic.
\end{theorem}

Taking $N=np^r$ gives $\vp(N)=r+\vp(n)\ge r$, explaining why the stated
results imply precisely the moduli requested in the OEIS entry, even
when $p\mid n$. The two identities in \eqref{eq:identity-uv} also show
that the apparently different conjectures belong to one normalized
family, with parameters differing by one.

The proof has three ingredients. First, a binomial-product identity
establishes order-two congruences for the seed $a_n$ at every prime.
Second, these congruences make $W$ a $2$-function and force an integral
Euler product for $G$. Finally, a change of variable preserves the relevant
Frobenius divisibility. Lagrange inversion identifies its coefficients
with $B_t(n)$.

\section{An elementary congruence for the Ap\'ery seed}\label{sec:seed}
It is known that Ap\'ery numbers satisfy the stronger order-three
congruences at primes at least $5$; the historical results of Gessel and
Coster are reviewed in \cite{liu}. We do not need that stronger theorem.
The following direct argument proves all the input used below, including
at $2$ and $3$.

\begin{lemma}[Binomial product]\label{lem:product}
For integers $1\le k\le N$, put
$C(N,k)=\binom Nk\binom{N+k}k$. Then
\begin{equation}\label{eq:product}
 C(N,k)=(-1)^{k-1}\frac{N(N+k)}{k^2}
       \prod_{j=1}^{k-1}\left(1-\frac{N^2}{j^2}\right).
\end{equation}
If $p\mid N$ and $1\le\ell\le N/p$, then
\begin{equation}\label{eq:split-product}
 C(N,p\ell)=(-1)^{(p-1)\ell}C(N/p,\ell)
 \prod_{\substack{1\le j<p\ell\\p\nmid j}}
       \left(1-\frac{N^2}{j^2}\right).
\end{equation}
\end{lemma}
\begin{proof}
Extract $N/k$ from $\binom Nk$ and $(N+k)/k$ from $\binom{N+k}k$.
Pairing the remaining factors gives
$(N-j)(N+j)/j^2=-(1-N^2/j^2)$ and proves \eqref{eq:product}.
When $k=p\ell$, the factors with $p\mid j$ form exactly the product in
\eqref{eq:product} for $(N/p,\ell)$. The prefactors coincide. The ratio
of signs is $(-1)^{p\ell-1-(\ell-1)}=(-1)^{(p-1)\ell}$, proving
\eqref{eq:split-product}.
\end{proof}

\begin{proposition}[Order two at every prime]\label{prop:seed}
For every prime $p$ and every positive $N$ divisible by $p$,
\begin{equation}\label{eq:seed}
 a_N\equiv a_{N/p}\pmod{p^{2\vp(N)}}.
\end{equation}
\end{proposition}
\begin{proof}
Write $r=\vp(N)$. If $p\nmid k$, the identity
$\binom Nk=(N/k)\binom{N-1}{k-1}$ gives
$\vp(\binom Nk)\ge r$. Hence $C(N,k)^2$ is divisible by $p^{2r}$.

If $k=p\ell$, each denominator $j$ in the remaining product of
\eqref{eq:split-product} is a $p$-adic unit. Thus each factor is
$1$ modulo $p^{2r}$. On squaring, the sign disappears, and we obtain
\[
 C(N,p\ell)^2\equiv C(N/p,\ell)^2\pmod{p^{2r}}.
\]
This holds also for $p=2$: no assertion that the unsquared sign is
positive has been made. The term $k=0$ is $1$ on both sides. Summing over
all $k$ proves \eqref{eq:seed}.
\end{proof}

\begin{remark}
One must split the product before reducing it modulo a prime power.
Factors with $p\mid j$ need not be congruent to $1$ and cannot be discarded.
Their exact contribution is the lower-index binomial product.
\end{remark}

\section{Dilogarithms, Euler products, and local defects}\label{sec:twofunctions}
We now isolate the general arithmetic structure. The subsequent argument
works for any integer sequence satisfying \eqref{eq:seed}, not just the
Ap\'ery numbers.

\begin{definition}
An integer sequence $(a_n)_{n\ge1}$ will be called \emph{order-two
admissible} if \eqref{eq:seed} holds for every prime and every positive
multiple $N$ of that prime. Its associated series $W$ in
\eqref{eq:basic-series} is then a rational $2$-function.
\end{definition}

\begin{proposition}[Three equivalent descriptions]\label{prop:equivalence}
For an integer sequence $(a_n)_{n\ge1}$, the following are equivalent:
\begin{enumerate}[label=\textup{(\roman*)},nosep]
\item The sequence is order-two admissible.
\item For every $d\ge1$,
\begin{equation}\label{eq:cd}
 c_d=\frac1{d^2}\sum_{e\mid d}\mu(e)a_{d/e}\in\Z.
\end{equation}
\item For every prime $p$,
\begin{equation}\label{eq:R}
 R_p(z):=\frac1{p^2}W(z^p)-W(z)\in z\Zp[[z]].
\end{equation}
\end{enumerate}
When these conditions hold,
\begin{align}
 W(z)&=\sum_{d\ge1}c_d\Li_2(z^d),\label{eq:dilog}\\
 G(z)&=\prod_{d\ge1}(1-z^d)^{-dc_d}\in1+z\Z[[z]].\label{eq:euler}
\end{align}
\end{proposition}
\begin{proof}
Suppose (i), fix $d=p^rh$ with $p\nmid h$, and group terms in the
numerator of \eqref{eq:cd} according to whether the squarefree divisor
$e$ contains $p$. The numerator becomes
\[
 \sum_{e\mid h}\mu(e)
       \bigl(a_{p^rh/e}-a_{p^{r-1}h/e}\bigr).
\]
Every summand is divisible by $p^{2r}$. Doing this for every prime dividing
$d$ proves divisibility by $d^2$, hence (ii).

M\"obius inversion gives
\begin{equation}\label{eq:divisor-sum}
 a_n=\sum_{d\mid n}d^2c_d.
\end{equation}
Conversely, if (ii) holds and $N=p^rh$ with $p\nmid h$, subtracting the
expressions \eqref{eq:divisor-sum} for $N$ and $N/p$ leaves only divisors
$d$ with $\vp(d)=r$. Their squares are multiples of $p^{2r}$, proving (i).

For (iii), the coefficient of $z^N$ in $R_p$ is
\begin{equation}\label{eq:defect-coeff}
 [z^N]R_p(z)=
 \begin{cases}
 -a_N/N^2,&p\nmid N,\\
 (a_{N/p}-a_N)/N^2,&p\mid N.
 \end{cases}
\end{equation}
Since the $a_N$ are integers, the equivalence of (i) and (iii) follows.
Finally, expand $\Li_2(z)=\sum_{j\ge1}z^j/j^2$ and apply
\eqref{eq:divisor-sum} to obtain \eqref{eq:dilog}. Because
$\delta_z\Li_2(z^d)=-d\log(1-z^d)$, exponentiation gives
\eqref{eq:euler}. All products and sums are locally finite as formal
series. Integer, including negative, powers of an integral unit have
integer coefficients, which proves the last assertion.
\end{proof}

For the Ap\'ery seed, $c_1,\ldots,c_6$ are
$5,17,160,2058,32760,596092$. For example, the first factors of $G$ are
\[
 (1-z)^{-5}(1-z^2)^{-34}(1-z^3)^{-480}(1-z^4)^{-8232}\cdots.
\]
Their integrality is a consequence of the congruences, not an assumption
based on the initial coefficient list.

The logarithmic derivative of \eqref{eq:R} is also integral:
\begin{equation}\label{eq:D}
 D_p(z):=\delta_zR_p(z)
 =\frac1pL(z^p)-L(z)\in z\Zp[[z]].
\end{equation}
Although $W$ and $L$ can have denominators, all higher derivatives needed
below are integral:
\begin{equation}\label{eq:derivatives}
 \delta_z^kW(z)=\sum_{n\ge1}n^{k-2}a_nz^n
 \in z\Z[[z]]\qquad(k\ge2).
\end{equation}
This distinction between the first derivative and the higher derivatives
is important in the cancellation argument.

\section{The normalized family as a change of variable}\label{sec:lagrange}
For an integer $t$, introduce
\begin{equation}\label{eq:coordinate}
 f_t(z)=zG(z)^{-t},\qquad
 \psi_t(w)=\Rev(f_t)(w).
\end{equation}
Since $G$ is an integral unit, $f_t(z)=z+O(z^2)$ has integer coefficients.
Its inverse $\psi_t(w)=w+O(w^2)$ has integer coefficients as well: the
coefficient of degree $n$ in the equation $f_t(\psi_t(w))=w$ determines
the new inverse coefficient with coefficient $1$, using only integer
polynomials in earlier coefficients.

Set
\begin{equation}\label{eq:framed-potential}
 W_t(w)=W(\psi_t(w))-\frac t2L(\psi_t(w))^2.
\end{equation}
The quadratic correction is essential. Merely composing $W$ with
$\psi_t$ does not yield the cancellation used later.

\begin{lemma}[Differential and coefficient identities]\label{lem:lagrange}
For every integer $t$,
\begin{align}
 \delta_wW_t(w)&=L(\psi_t(w)),\label{eq:framed-derivative}\\
 W_t(w)&=\sum_{n\ge1}\frac{B_t(n)}{n^2}w^n.\label{eq:framed-coefficients}
\end{align}
\end{lemma}
\begin{proof}
In the coordinate $w=zG(z)^{-t}$,
\[
 \frac{\dd w}{w}=(1-t\delta_zL(z))\frac{\dd z}{z},
 \qquad
 \dd\left(W(z)-\frac t2L(z)^2\right)
 =L(z)(1-t\delta_zL(z))\frac{\dd z}{z}.
\]
The factor $1-t\delta_zL(z)$ has constant term $1$, so is a formal unit.
Division gives \eqref{eq:framed-derivative}.

The Lagrange coefficient formula, proved in \cref{app:lagrange}, says that
if $\psi(w)=w\Phi(\psi(w))$, then for $n\ge1$,
\begin{equation}\label{eq:lagrange}
 [w^n]H(\psi(w))=\frac1n[z^{n-1}]H'(z)\Phi(z)^n.
\end{equation}
Here $\Phi=G^t$ and $H=L$. When $t\ne0$,
\begin{align*}
 [w^n]L(\psi_t(w))
 &=\frac1n[z^{n-1}]L'(z)G(z)^{tn}\\
 &=\frac1{tn^2}[z^{n-1}]\bigl(G(z)^{tn}\bigr)'\\
 &=\frac1{tn}[z^n]G(z)^{tn}=\frac{B_t(n)}n.
\end{align*}
For $t=0$, $\psi_0(w)=w$ and the same identity follows from
$B_0(n)=a_n$. Both sides of \eqref{eq:framed-coefficients} have zero
constant coefficient, so logarithmic integration proves it.
\end{proof}

At this point $B_t(n)$ is only known to be rational when $t\ne0$.
The next section proves the needed local integrality before using it to
assert divisibility. This avoids treating division by $t$ as though it
were harmless at a prime dividing $t$.

\section{The Frobenius-defect calculation}\label{sec:framing}
This section gives the central proof. It uses only the conclusions of
\cref{prop:equivalence}, the integral coordinate inverse, and elementary
bounds on factorial denominators.

\subsection{An integral logarithmic error}
Fix a prime $p$ and an integer $t$. Write
\begin{equation}\label{eq:x-and-J}
 x=\psi_t(w),\qquad x_p=\psi_t(w^p),\qquad
 J(w)=G(\psi_t(w))\in1+w\Z[[w]].
\end{equation}
Define
\begin{equation}\label{eq:S}
 S(w)=\frac1p\log J(w^p)-\log J(w)
     =\frac1pL(x_p)-L(x).
\end{equation}

\begin{lemma}\label{lem:S}
For every prime $p$, including $2$,
$S(w)\in w\Zp[[w]]$. Moreover,
\begin{equation}\label{eq:xp}
 x_p=x^p\exp\bigl(ptS(w)\bigr).
\end{equation}
\end{lemma}
\begin{proof}
Reduction modulo $p$ gives $J(w^p)\equiv J(w)^p$ because the coefficients
are integers. Thus
\[
 \frac{J(w^p)}{J(w)^p}=1+pH(w),\qquad H(w)\in w\Zp[[w]].
\]
The formal logarithm yields
\[
 S(w)=\frac1p\log(1+pH(w))
     =\sum_{j\ge1}\frac{(-1)^{j+1}p^{j-1}}jH(w)^j.
\]
The scalar coefficients are $p$-integral, since
$j-1-\vp(j)\ge0$ for every $j\ge1$. This includes $p=2$.
For example, if $e=\vp(j)$ then $j\ge2^e\ge e+1$.
The sum is well defined coefficientwise because $H$ has zero constant term.

The inverse-coordinate equation says $x=wJ(w)^t$. Applying it at $w^p$
and dividing by $x^p=w^pJ(w)^{pt}$ gives
\[
 \frac{x_p}{x^p}
 =\left(\frac{J(w^p)}{J(w)^p}\right)^t
 =\exp(ptS(w)),
\]
which is \eqref{eq:xp}. All quotients here are quotients of formal units
after the common leading monomial has been canceled.
\end{proof}

\subsection{The exact cancellation identity}
The framed Frobenius defect is
\begin{equation}\label{eq:Rtilde}
 \widetilde R_{p,t}(w)=\frac1{p^2}W_t(w^p)-W_t(w).
\end{equation}
The goal is to prove its coefficients integral, allowing at most one
factor of $2$ in their denominators when both $p=2$ and $t$ is odd.

\begin{proposition}[Explicit defect formula]\label{prop:defect-identity}
With the notation above,
\begin{equation}\label{eq:key}
\boxed{\begin{aligned}
 \widetilde R_{p,t}(w)
 &=R_p(x)+tS(w)D_p(x)-\frac t2S(w)^2\\
 &\quad+\sum_{k\ge2}\frac{p^{k-2}t^k}{k!}S(w)^k
                  (\delta_z^kW)(x^p).
\end{aligned}}
\end{equation}
This is an equality of formal series over $\Qp$.
\end{proposition}
\begin{proof}
For a series $V$ with zero constant coefficient, the logarithmic Taylor
identity is
\begin{equation}\label{eq:logtaylor}
 V(ye^q)=\sum_{k\ge0}\frac{q^k}{k!}(\delta_y^kV)(y).
\end{equation}
It follows by checking the monomials $y^n$: both sides are $y^ne^{nq}$.
Afterward one may substitute series for $y$ and $q$ having zero constant
term. Every coefficient then involves only finitely many terms.

Apply \eqref{eq:logtaylor} to $W$, using $y=x^p$ and $q=ptS$:
\begin{align*}
 \frac1{p^2}W(x_p)
 &=\frac1{p^2}W(x^p)+\frac{tS}{p}L(x^p)\\
 &\quad+\sum_{k\ge2}\frac{p^{k-2}t^k}{k!}S^k
                          (\delta_z^kW)(x^p).
\end{align*}
Also, by \eqref{eq:S}, $L(x_p)/p=L(x)+S$. The quadratic part of
\eqref{eq:framed-potential} therefore contributes
\[
 -\frac t2\left(\left(\frac{L(x_p)}p\right)^2-L(x)^2\right)
 =-tS L(x)-\frac t2S^2.
\]
Combining the linear Taylor term with the first term of this expression
gives
\[
 tS\left(\frac{L(x^p)}p-L(x)\right)=tS D_p(x).
\]
The terms involving $W$ with no Taylor derivatives give $R_p(x)$.
These substitutions prove \eqref{eq:key}.
\end{proof}

The terms $L(x)$ and $L(x^p)/p$ need not be integral separately.
Their cancellation is what prevents the denominator $n$ in the logarithm
from costing a power of the modulus. After this cancellation, only
integral series and explicit factorial denominators remain.

\subsection{Factorial denominators and the prime two}
\begin{lemma}\label{lem:factorial}
For $k\ge2$ and an odd prime $p$,
\begin{equation}\label{eq:factorial-odd}
 \vp(k!)\le k-2.
\end{equation}
For $p=2$,
\begin{equation}\label{eq:factorial-two}
 \vtwo(k!)\le k-1.
\end{equation}
Consequently $p^{k-2}/k!$ is $p$-integral for odd $p$, and
$2^{k-2}/k!\in\tfrac12\Z_2$.
\end{lemma}
\begin{proof}
Counting multiples of each power of $p$ among $1,\ldots,k$ gives
$\vp(k!)=\sum_{j\ge1}\lfloor k/p^j\rfloor$.
Equivalently, if $s_p(k)\ge1$ is the sum of the base-$p$ digits of $k$,
this is $(k-s_p(k))/(p-1)$. For odd $p$ and $k\ge3$, it is at most
$(k-1)/2\le k-2$; $k=2$ is immediate. At $p=2$ it is at most $k-1$.
\end{proof}

\begin{theorem}[Unsigned framing, with the $2$-adic loss]\label{thm:framing}
Let $(a_n)_{n\ge1}$ be any order-two admissible integer sequence,
and define $L,G,W,\psi_t,W_t$ as above. For every integer $t$,
\begin{equation}\label{eq:defect-integrality}
 \widetilde R_{p,t}(w)\in
 \begin{cases}
 w\Zp[[w]],&p\text{ odd},\\
 w\Z_2[[w]],&p=2\text{ and }t\text{ even},\\
 \tfrac12 w\Z_2[[w]],&p=2\text{ and }t\text{ odd}.
 \end{cases}
\end{equation}
Its normalized coefficients $B_t(n)$ are integers for all integers $t$,
and they satisfy \eqref{eq:B-odd} and \eqref{eq:B-two}.
\end{theorem}
\begin{proof}
By \cref{prop:equivalence,lem:S}, the series $R_p(x)$, $D_p(x)$, and $S$
are $p$-integral. By \eqref{eq:derivatives}, so is
$(\delta_z^kW)(x^p)$ for every $k\ge2$. Composition preserves integrality
because $x\in w\Z[[w]]$.

For odd $p$, the number $1/2$ is $p$-integral, and
\cref{lem:factorial} makes every scalar $p^{k-2}t^k/k!$ integral.
Thus \eqref{eq:key} lies in $w\Zp[[w]]$.

For $p=2$, the same argument allows at worst a factor $1/2$, both from
$tS^2/2$ and from the factorial coefficients. If $t$ is even, the former
term is integral and so are the latter coefficients: for nonzero even
$t$ and $k\ge2$,
\[
 \vtwo\left(\frac{2^{k-2}t^k}{k!}\right)
 \ge k-2+k-(k-1)=k-1\ge1.
\]
For $t=0$ all those terms vanish. This proves
\eqref{eq:defect-integrality}.

We now deduce the coefficient assertions without assuming in advance
that $B_t(n)$ is integral. From \eqref{eq:framed-coefficients},
\begin{equation}\label{eq:framed-defect-coeff}
 [w^N]\widetilde R_{p,t}(w)=
 \begin{cases}
 -B_t(N)/N^2,&p\nmid N,\\
 (B_t(N/p)-B_t(N))/N^2,&p\mid N.
 \end{cases}
\end{equation}
When the defect is integral, the first case shows that $B_t(N)$ is
$p$-integral for $p\nmid N$. The second case then proves, by induction
on $\vp(N)$, that every $B_t(N)$ is $p$-integral; simultaneously it proves
that $B_t(N)-B_t(N/p)$ has valuation at least $2\vp(N)$.
This covers every odd prime and also $p=2$ when $t$ is even.

For odd $t$, the definition \eqref{eq:Bt} has numerator in $\Z$, since
$G^{tn}$ is an integer power of an integral unit, and denominator $t$
is odd. Hence $B_t(n)$ is $2$-integral in that case as well. The half-integral
defect in \eqref{eq:framed-defect-coeff} gives exactly the lower bound
$2\vtwo(N)-1$ for the difference. The case $t=0$ is also immediate from
$B_0(n)=a_n$. Thus every rational $B_t(n)$ is integral at every prime,
and therefore belongs to $\Z$. All the asserted bounds follow.
\end{proof}

\Cref{prop:seed} supplies the hypotheses for Ap\'ery numbers, so
\cref{thm:framing} proves \cref{thm:normalized}. The proof makes no
coprimality assumption between $t$ and $p$. Its only use of the
ordinary quotient by $t$ for integrality is at $p=2$ when $t$ is odd,
where that quotient is legitimate.

\section{Finishing both conjectures}\label{sec:finish}
The first identity in \eqref{eq:identity-uv} is immediate for $m\ne0$,
and for $m=0$ both sides are zero at positive indices. The reversion
identity requires a separate coefficient calculation.

\begin{proposition}[A parameter shift accounts for reversion]\label{prop:shift}
For every integer $m$ and every $n\ge1$,
\[
 [w^n]F(w)^{mn}=mB_{m-1}(n).
\]
In particular, $[w^n]F(w)^n=a_n$.
\end{proposition}
\begin{proof}
The inverse $h(w)$ in \eqref{eq:F} is $\psi_{-1}(w)$.
The equation $w=h(w)G(h(w))$ gives
\[
 F(w)=G(h(w))\in1+w\Z[[w]].
\]
Apply \eqref{eq:lagrange} to $h=wG(h)^{-1}$ and
$H(z)=G(z)^{mn}$. Since integer powers, including negative powers, are
well defined formally,
\begin{align*}
 [w^n]F(w)^{mn}
 &=\frac1n[z^{n-1}]\bigl(G(z)^{mn}\bigr)'G(z)^{-n}\\
 &=m[z^{n-1}]L'(z)G(z)^{(m-1)n}.
\end{align*}
For $m\ne1$, differentiating $G^{(m-1)n}$ turns the last expression into
\[
 \frac{m}{m-1}[z^n]G(z)^{(m-1)n}=mB_{m-1}(n).
\]
When $m=1$, the previous line before division is
$[z^{n-1}]L'(z)=a_n=B_0(n)$. The case $m=0$ gives zero directly.
These are equalities over $\Q$, so no modular division by $m-1$ is used.
\end{proof}

\begin{proof}[Proof of \cref{thm:OEIS}]
The exact identities \eqref{eq:identity-uv} have now been established.
For odd $p$, multiply \eqref{eq:B-odd} by $m$ to obtain
\eqref{eq:strong-u} and \eqref{eq:strong-v} for $m\ne0$. They imply the
original congruences, and the same proof applies at $p=3$.

For $u_m$ at $2$, the framing parameter is $t=m$. If $m$ is even,
\eqref{eq:B-two} gives full precision $2\vtwo(N)$, and multiplication
by $m$ supplies $\vtwo(m)$ more powers. If $m$ is odd, it gives precision
$2\vtwo(N)-1$ and multiplication loses or gains no power of $2$.
This is \eqref{eq:u-two}.

For $v_m$ the framing parameter is $m-1$. If $m$ is odd, that parameter
is even, so \eqref{eq:B-two} gives full precision. If $m$ is nonzero
even, the normalized family can lose one power of $2$, but the prefactor
$m$ supplies at least one. More precisely the resulting bound is
$2\vtwo(N)+\vtwo(m)-1$. Thus \eqref{eq:v-two} holds, and its lower bound
is at least $2\vtwo(N)$ in both cases. This proves the second OEIS
congruence at $2$ for every $m$, including the trivial case $m=0$.
\end{proof}

The value $v_1(n)=a_n$ recovers the familiar sequence rather than creating
an unrelated one. Another useful exact relation is
\begin{equation}\label{eq:v2}
 v_2(n)=2B_1(n)=2u_1(n).
\end{equation}
More generally, if $m\ne1$,
$v_m(n)=m\,u_{m-1}(n)/(m-1)$ as a rational identity. The normalized
form is preferable for divisibility arguments because the integrality of
$B_{m-1}(n)$ has already been proved even at primes dividing $m-1$.

\section{Examples and limits of the result}\label{sec:examples}
\subsection{Integer-valued parameter polynomials}
The first few normalized polynomials are
\begin{align*}
 B_t(1)&=5,\\
 B_t(2)&=50t+73,\\
 B_t(3)&=\frac{1125t^2+3285t+2890}{2},\\
 B_t(4)&=\frac{20000t^3+87600t^2+147574t+99003}{3}.
\end{align*}
Their coefficients need not be integers. The theorem asserts that their
values at every integer $t$ are integers. This is a genuine
integer-valued-polynomial conclusion, not a statement that all coefficient
denominators disappear algebraically.

For completeness, expanding the exponential gives the finite expression
\begin{equation}\label{eq:partitionformula}
 B_t(n)=\sum_{\substack{k_1,\ldots,k_n\ge0\\
                        k_1+2k_2+\cdots+nk_n=n}}
 n^K t^{K-1}\prod_{j=1}^n\frac{(a_j/j)^{k_j}}{k_j!},
 \qquad K=k_1+\cdots+k_n.
\end{equation}
Since $n\ge1$, every summand has $K\ge1$. At $t=0$ only $K=1$ survives,
and that term is $a_n$. Formula \eqref{eq:partitionformula} is useful
for symbolic experimentation, but its denominators do not by themselves
explain either integrality or the supercongruences.

\begin{table}[htbp]
\centering
\small
\setlength{\tabcolsep}{9pt}
\begin{tabular}{rrrrrr}
\toprule
$t$ & $B_t(1)$ & $B_t(2)$ & $B_t(3)$ & $B_t(4)$ & $B_t(5)$\\
\midrule
$-2$ & 5 & $-27$ & 410 & $-1915$ & 53130\\
$-1$ & 5 & 23 & 365 & 6343 & 123255\\
0 & 5 & 73 & 1445 & 33001 & 819005\\
1 & 5 & 123 & 3650 & 118059 & 4015380\\
2 & 5 & 173 & 6980 & 301517 & 13540505\\
3 & 5 & 223 & 11435 & 623375 & 35175630\\
\bottomrule
\end{tabular}
\caption{Exact values, including negative framing parameters.}\label{tab:values}
\end{table}

\subsection{The obstruction at two is real}
For $m=1$,
\[
 u_1(1)=5,\qquad u_1(2)=[z^2]G(z)^2=2\cdot49+5^2=123.
\]
Therefore
\begin{equation}\label{eq:2counter}
 u_1(2)-u_1(1)=118\equiv2\pmod4.
\end{equation}
So an assertion that the first conjecture holds for all primes and all
integers $m$ would be false. The lower bound
$2\vtwo(2)-1=1$ in \eqref{eq:u-two} is attained.
In contrast, $v_2(2)-v_2(1)=2\cdot118=236$ is divisible by $4$, illustrating
the compensating factor in the second family. This does not claim that
every individual odd parameter always fails the stronger $2$-adic modulus;
it disproves the uniform assertion.

\subsection{A uniform cubic modulus is false}
The initial Ap\'ery numbers have order-three congruences for $p\ge5$,
but a general framing need not preserve that order.
Exact coefficient extraction gives
\[
 u_1(7)=5\,082\,313\,276,
\]
and hence
\begin{align*}
 u_1(7)-u_1(1)
 &=5\,082\,313\,271\\
 &=7^2\cdot103\,720\,679,
 \qquad 103\,720\,679\equiv6\pmod7.
\end{align*}
Thus its $7$-adic valuation is exactly $2$, not $3$, and its residue
modulo $343$ is $294$. By \eqref{eq:v2},
\[
 v_2(7)-v_2(1)=10\,164\,626\,542\equiv245\pmod{343},
\]
again with valuation exactly $2$. Therefore neither conjecture admits a
uniform replacement of $p^{2r}$ by $p^{3r}$ for all parameters, even in the
original range $p\ge5$. These are counterexamples to stronger variants,
not to either formula actually listed in OEIS.

\section{A broader family covered by the same argument}\label{sec:generalization}
The transform theorem is not specific to one recurrence. It applies to
any integer seed satisfying \eqref{eq:seed}. There is also a direct
extension of the elementary binomial argument.

\begin{proposition}\label{prop:moments}
For every positive integer $h$, the binomial-moment sequence
\[
 a_n^{(h)}=\sum_{k=0}^n
       \left(\binom nk\binom{n+k}k\right)^{2h}
\]
is order-two admissible. Consequently its associated normalized family
satisfies \cref{thm:normalized}, and its exponential and reverted families
satisfy \cref{thm:OEIS} with this seed in place of $a_n$.
\end{proposition}
\begin{proof}
If $p\nmid k$, the argument of \cref{prop:seed} gives
$\vp(C(N,k)^{2h})\ge2h\vp(N)\ge2\vp(N)$.
If $k=p\ell$, raise \eqref{eq:split-product} to the even power $2h$.
The sign disappears and each remaining factor is $1$ modulo
$p^{2\vp(N)}$. Summing proves the required seed congruence, after which
\cref{thm:framing,prop:shift} apply without change.
\end{proof}

The case $h=1$ is exactly the classical Ap\'ery sequence used to resolve
the two selected conjectures. This extension is included to show what
the proof depends on: an arithmetic property of a binomial sum and an
integral change of variable, not a numerical fit or an unproved
classification of recurrences.

\section{Exact computation and reproducibility}\label{sec:computations}
The accompanying program \texttt{verify.py} uses only the Python standard
library. It uses arbitrary-precision integers and \texttt{Fraction}; no
floating-point equality or probabilistic primality decision enters the
checks. The following describes the saved run with limit $150$.

\subsection{Two ways to generate the seed}
For efficiency, the program generates the seed by the classical recurrence
\cite{liu}
\begin{equation}\label{eq:recurrence}
 (n+1)^3a_{n+1}
 =(2n+1)(17n^2+17n+5)a_n-n^3a_{n-1},
 \qquad a_0=1,\ a_1=5.
\end{equation}
It independently evaluates the defining binomial sum at every index
$0\le n\le150$ and checks equality. Every division in the recurrence is
checked for zero remainder. The mathematical proof in this article uses
the binomial definition, not the recurrence, so no separate proof of
\eqref{eq:recurrence} is an unmentioned dependency of the main theorem.

\subsection{Exact coefficients and independent reversion}
For a fixed integer exponent $q$, write $G(z)^q=\sum_{j\ge0}h_jz^j$.
Logarithmic differentiation gives
\begin{equation}\label{eq:coefficientrec}
 h_0=1,\qquad jh_j=q\sum_{k=1}^j a_kh_{j-k}\quad(j\ge1).
\end{equation}
To compute $B_t(n)$, the program sets $q=tn$, evaluates the recurrence
through degree $n$, and divides by $t$ for $t\ne0$. All remainders are
checked; for $t=0$ the seed value is used.

Independently, the program constructs the compositional inverse of $zG(z)$
by solving its coefficients successively, checks the composition, forms
$F=w/h(w)$, and computes the powers $F^{mn}$ by rational polynomial
convolution. Those tests do not use $v_m(n)=mB_{m-1}(n)$ to generate the
left-hand side. They therefore provide a separate check of the parameter
shift and the reversion convention.

Finally, the full two sides of \eqref{eq:key} are constructed as truncated
rational series through degree $14$. This checks the quadratic sign, the
factor $p^{k-2}$, and the placement of $x^p$ as distinct from $x_p$, rather
than only comparing a selection of final congruences.

\begin{table}[htbp]
\centering
\small
\begin{tabular}{p{0.76\textwidth}r}
\toprule
Check performed in the saved run & Count\\
\midrule
Binomial-sum versus recurrence values, indices $0$--$150$ & 151\\
Seed congruences, all primes and multiples $N\le150$ & 269\\
Odd-prime normalized congruences on the specified grid & 798\\
Normalized $2$-adic bounds, all even $N\le150$ & 1575\\
Parameter-sensitive odd-prime bounds for $u_m$ & 840\\
Parameter-sensitive $2$-adic bounds for $u_m$ & 1575\\
Bounds for $v_m$, including the prime $2$ & 2415\\
Independent reversion-and-power identities & 126\\
Full Frobenius-defect identities through degree $14$ & 18\\
\bottomrule
\end{tabular}
\caption{All listed exact checks passed. These finite counts are validation,
not a substitute for the general proof.}\label{tab:checks}
\end{table}

The normalized odd-prime grid uses $t=-10,-9,\ldots,10$,
$p\in\{3,5,7,11,13\}$, $1\le r\le4$, and $1\le n\le4$, retaining
$N=np^r\le150$. The $2$-adic normalized tests use all even $N\le150$
for the same $t$ range. The parameter-sensitive $u_m$ and odd-prime
$v_m$ tests use $m=-10,\ldots,10$ and multiples $N\le50$ of those odd
primes; the $u_m$ and $v_m$ tests additionally check their full
parameter-sensitive bounds at every even $N\le150$.
The independent reversion tests use $m=-4,\ldots,4$ and $1\le n\le14$.
The $18$ formal defect tests use
$t\in\{-3,-1,0,1,2,4\}$ and $p\in\{2,3,5\}$.
The sharpness examples in \cref{sec:examples} are also checked exactly.

\subsection{Rebuilding the artifacts}
From the archive directory, the verification command is
\begin{verbatim}
python3 verify.py --out data --limit 150
\end{verbatim}
Run Python normally, not with the \texttt{-O} option: the verification
uses assertions. The output includes a JSON summary and CSV tables.
The PDF can be rebuilt without a separate bibliography processor:
\begin{verbatim}
pdflatex -interaction=nonstopmode -halt-on-error article.tex
pdflatex -interaction=nonstopmode -halt-on-error article.tex
\end{verbatim}
The README records the artifact inventory. Numerical checks cannot prove
an infinite family of congruences. Their role here is to make the examples
reproducible and to expose implementation, normalization, or sign errors
in independently computed quantities. The proofs establish the claims
for every parameter and index, beyond the finite test range.

\section{Conclusion and attribution}\label{sec:conclusion}
Both October 2024 supercongruence conjectures in A267220 are true. The
same argument strengthens the first from $p\ge5$ to every odd prime,
and the second to every prime. It proves integrality of the normalized
family, gains powers from the integer parameter, and identifies the
precise general one-power loss at $2$ for odd unsigned framings.
Counterexamples show why two natural stronger uniform assertions must
not be inferred.

The structural explanation is that the Ap\'ery binomial sum produces a
$2$-function, and the two OEIS constructions read off coefficients after
closely related changes of variable. The known framing theory of
\cite{svw2013,svw2017} is therefore the appropriate context. The article
has not asserted that this general principle, the classical Ap\'ery
congruences, or the old integrality remarks in OEIS are new. Nor does
an entry's conjectural label establish that no proof exists elsewhere.
What is supplied here is an explicit self-contained derivation of the
two listed claims, with the prime-$2$ refinement and all integer-parameter
edge cases accounted for.

\appendix
\section{A formal proof of the Lagrange formula}\label{app:lagrange}
We include the coefficient formula to make the power-series part of the
argument self-contained. Let $\Phi(z)\in1+z\Q[[z]]$, put
$f(z)=z/\Phi(z)$, and let $\psi$ be its inverse. For a Laurent series,
$\operatorname{Res}_z$ means the coefficient of $z^{-1}$ in a formal
one-form. Formal residue is invariant under an invertible change of
variable, so
\begin{align*}
 [w^n]H(\psi(w))
 &=\operatorname{Res}_w H(\psi(w))w^{-n-1}\dd w\\
 &=\operatorname{Res}_z H(z)f(z)^{-n-1}f'(z)\dd z.
\end{align*}
For clarity, residue invariance can be checked on each monomial $w^j$:
if $j\ne-1$, the substituted form is the derivative of
$f(z)^{j+1}/(j+1)$ and has zero residue; if $j=-1$, $f'/f$ has residue $1$
because $f=z$ times a unit.

Using $\dd(f^{-n})=-n f^{-n-1}f'\dd z$ and the fact that a formal derivative
has zero residue, integration by parts gives
\begin{align*}
 \operatorname{Res}_z H(z)f(z)^{-n-1}f'(z)\dd z
 &=\frac1n\operatorname{Res}_z H'(z)f(z)^{-n}\dd z\\
 &=\frac1n[z^{n-1}]H'(z)\Phi(z)^n.
\end{align*}
This is \eqref{eq:lagrange}, for all positive $n$. The constant term of
$H$ is irrelevant to both sides.

\section{Proof dependencies and boundary cases}\label{app:audit}
For an independent audit, the logical dependencies can be summarized
without treating computation as a lemma. The binomial identity proves
\cref{prop:seed}. M\"obius inversion proves \cref{prop:equivalence}, including
integrality of $G$. Integral compositional inversion and the Lagrange
formula identify the coefficients of $W_t$. The elementary logarithmic
lemma, exact Taylor cancellation, and factorial bounds prove
\cref{thm:framing}. The parameter-shift identity then proves both OEIS
statements.

The points most vulnerable to an incorrect shortcut are the following.
The inverse is \emph{compositional}, and $F=G\circ\psi_{-1}$, not
$G\circ\psi_1$. The quadratic correction has sign $-t/2$ for the coordinate
$w=zG^{-t}$. The Frobenius Taylor expansion uses $x^p$, whereas $x_p$
denotes $\psi_t(w^p)$; these are generally different. The normalized
coefficient $B_t(n)$ is not assumed integral before its local proof.
Division by $m-1$ is avoided at $m=1$ and never used modulo $p$.
Negative parameters are integer powers of units and need no positivity
assumption. The two-adic half-integrality is retained rather than silently
replaced by integrality. Finally, the modulus depends on the full
valuation of $N$, not on a representation of $N$ with a tacitly
prime-to-$p$ leading factor.

At index zero the original definitions give $u_m(0)=v_m(0)=1$, even for
$m=0$. All normalized identities in this article are stated for positive
indices. The conjectures themselves have positive $n$ and $r$, so this
convention causes no exceptional case in their application.

\begin{thebibliography}{9}
\bibitem{oeis-apery}
The OEIS Foundation Inc.,
\emph{A005259: Ap\'ery numbers},
\url{https://oeis.org/A005259}.
Accessed 20 September 2026.

\bibitem{oeis-target}
The OEIS Foundation Inc.,
\emph{A267220: Expansion of the exponential associated with A005259},
\url{https://oeis.org/A267220}.
In particular, Peter Bala's two conjectures in the formula section,
dated 17 October 2024. Accessed 20 September 2026.

\bibitem{svw2013}
A. Schwarz, V. Vologodsky, and J. Walcher,
\emph{Framing the Di-Logarithm (over $\Z$)},
arXiv:1306.4298, 2013,
\url{https://arxiv.org/abs/1306.4298}.

\bibitem{svw2017}
A. Schwarz, V. Vologodsky, and J. Walcher,
\emph{Integrality of Framing and Geometric Origin of 2-functions
(with algebraic coefficients)},
arXiv:1702.07135v2, 6 March 2017,
\url{https://arxiv.org/abs/1702.07135}.

\bibitem{liu}
J.-C. Liu,
\emph{An extension of Gauss congruences for Ap\'ery numbers},
arXiv:2404.16636v1, 25 April 2024,
\url{https://arxiv.org/abs/2404.16636}.
Used for background, the standard recurrence, and the historical
order-three congruences; the input order-two congruence is proved directly
in this article.
\end{thebibliography}
\end{document}
